\section{Causal Discovery \& Estimation with iSCMs}

So far, we always assumed that an iSCM was fully specified, and derived theory to 
draw conclusions from the given iSCM. For example, the do-calculus provides precise
relationships between certain Markov kernels induced by the iSCM,
which enables us to perform \emph{causal reasoning}.

However, often we do not have sufficient information regarding the system that
we are modeling to completely specify an iSCM. For example, we may only know what
the observed variables are, but not what the graph of the iSCM is, let alone know the
latent spaces, exogenous distribution and exact causal mechanisms. Can we still
perform causal reasoning with such incompletely specified models? The answer turns
out to be affirmative, if one is willing to make certain assumptions (that---unfortunately---are 
typically untestable, even when having access to the observable Markov kernel of the model).

In the rest of this chapter we will focus on the question of how to deduce partial knowledge about
the iSCM from given Markov kernels. When this knowledge pertains only to the causal graph of the iSCM,
this is often called \emph{causal discovery}. When it pertains to parameters of the iSCM, for example,
when estimating the causal effect of a variable on another, this is often referred to as \emph{causal 
inference} (although ``inference'' is often interpreted much broader as drawing conclusions from data 
and prior beliefs).

In the next chapter, we will go one step further, and replace the deduction
of iSCM properties from Markov kernels by the estimation of iSCM properties
from data, i.e., we replace Markov kernels by finite samples. This will lead to 
statistical considerations. In particular, we will focus on estimating the
conditional independences of an iSCM from data.

In this chapter, we will make use of the simple iSCM formalism. 
Similar (more restrictive) results can be obtained in the L-CBN formalism.

\subsection{Randomized Controlled Trials}

The notion of randomized controlled trials (also known as A/B-testing in engineering), 
is centuries old. It was already proposed by the Flemish physician Jan Baptista van
Helmont \cite{vanHelmont1648} in 1648. As of today, it still provides the `gold standard'
for discovering causal relations and for the estimation of causal effects.

The experimental procedure is as follows.
Consider two variables, ``treatment'' $C$ and ``outcome'' $X$. 
In the simplest setting, one considers a 
binary treatment variable, where $C=1$ corresponds to ``treat with drug'' and
$C=0$ corresponds to ``treat with placebo'' in a medical setting, or with ``arm A''
and ``arm B'' in an engineering setting. For example, the drug could be aspirin,
and outcome could be the severity of headache perceived two
hours later. Patients are split into two groups,
the treatment and the control group, by means of a coin flip that assigns a value
of $C$ to every patient.\footnote{Usually this is done in a double-blind way, so
that neither the patient nor the doctor knows which group a patient has been assigned
to.} Patients are treated depending on the assigned value of $C$, i.e., patients in the
treatment group are treated with the drug and patients in the control group are treated
with a placebo. Some time after treatment, the outcome $X$ is measured for each patient. 
This yields a data set $(C_n,X_n)_{n=1}^N$
with two measurements ($C_n$, $X_n$) for the $n^{\mathrm{th}}$ patient. 
If the distribution of outcome $X$ significantly differs between the two groups,
one concludes that treatment is a cause of outcome. 

Let us formalize this in the causal modeling language of iSCMs.
%\footnote{From now on, we will mostly work with iSCMs, although if feedback cycles can be ruled out a priori, one could also use L-CBNs instead.}
Apart from that treatment may have a causal 
effect on outcome, there are likely many other factors that influence outcome. Some
have been measured, others not. For obvious practical reasons, we are not going 
to explicitly model \emph{all} of them. Formally, we will assume that an accurate
causal model of the situation is provided by some (unknown) simple iSCM with observed
variables $C$ and $X$, and possibly other latent variables.
We will consider the outcome variable $X$ as endogenous. 
But what type of variable should we consider
the treatment variable $C$ to be (which is not necessarily binary)?
We have two possibilities: exogenous input, or endogenous. We will
discuss both of these possibilities in sequence.

Let us start by considering the treatment variable $C$ as an exogenous input variable. 
This choice encodes into the iSCM that outcome $X$ does not cause treatment $C$.
We are interested in answering two questions. The first is
``Does treatment cause outcome according to $G(M)$?'', where $M$ is the underlying (unknown) iSCM. 
Since there are only two observed variables, this is equivalent to asking
``Is $C \tuh X$ in $G(M)$?''.
The second question is ``What is the causal effect of treatment on outcome?''.
We interpret this as asking for the Markov kernel $P_{M}(X \mid \doit(C))$.
\begin{Prp}\label{prp:RCT1}
Let $M$ be a simple iSCM with a single exogenous input variable $C$ and a single endogenous variable $X$.
A dependence 
  \begin{equation}\label{eq:RCT_CI_statement1}
    X \nIndep_{P_{M}(X\mid\doit(C))} C 
  \end{equation}
implies that $C$ causes $X$ according to $G(M)$. 
\end{Prp}
\begin{proof}
  This follows immediately from Theorem~\ref{thm:scm_causes}.
\end{proof}
The second question is, in this case, trivially answered.

\begin{comment}
Alternatively, we can consider the treatment variable as an exogenous \emph{random} variable.
\begin{Prp}\label{prp:RCT2}
Let $\bar{M}$ be a simple iSCM with an exogenous random variable $C$ and an endogenous 
variable $X$, both of which are observed (and possibly other latent variables as well).
A dependence 
  \begin{equation}\label{eq:RCT_CI_statement2}
    X \nIndep_{P_{\bar{M}}(X,C)} C 
  \end{equation}
implies that $C$ causes $X$ according to $\bar{M}$.
The causal effect of $C$ on $X$ satisfies:
  \begin{equation}\label{eq:RCT_causal_effect2}
    P_{\bar{M}}\big(X \given \doit(C)\big) = P_{\bar{M}}(X \given C) \qquad P_{\bar{M}}(C)\text{-a.s.}.
  \end{equation}
\end{Prp}
\begin{proof}
Denote the observable graph as $\bar{G} := G_{\{C,X\}}(\bar{M})$.
It has two nodes, and it either has no edge at all, in which case $C$ does not cause $X$ according to
$\bar{M}$, or it has a single edge $C \tuh X$, in which case $C$ causes $X$ according to $\bar{M}$.
By the Markov property (Corollary~\ref{cor:gmp_simple_scm}), if the edge $C \tuh X$ were absent in $\bar{G}$, 
then $X \Indep_{P_{\bar{M}}(X,C)} C$. 
%Therefore, if $X \nIndep C$, the edge $C \to X$ must be in $\tilde{G}$. 
In both cases, rule 2 of the causal do-calculus applied to $\bar{G}$ yields the identity \eref{eq:RCT_causal_effect2}. 
\end{proof}
\end{comment}

The other option is to consider the treatment variable as \emph{endogenous}.
One situation in which this is more appropriate is so-called ``imperfect compliance''.
If trial subjects do not all comply with prescribed treatment, for whatever reasons, 
then we can no longer identify the coin flip outcome with treatment,
(even though coin flip outcome may still be an important cause of treatment). 

In this modeling variant, we assume the existence of a simple iSCM $M$
with endogenous variables $C, X$, both of which are observed (and possibly other
latent variables). $C$ here retains the meaning
of treatment (but is no longer necessarily identifiable with the coin
flip result).
%Again, we can invoke
%Corollary \ref{cor:scm_identify_causal_relation} to answer the question
%``Does treatment cause outcome?''. 
Under additional assumptions regarding the exogeneity of the treatment variable,
we again obtain a similar statement as before.
\begin{Prp}\label{prp:RCT3}
Let $\tilde{M}$ be a simple SCM with two observed endogenous variables $C, X$,
and no exogenous input variables.
Under the following two assumptions:
  \begin{enumerate}
    \item $X$ does not cause $C$ according to $G(\tilde{M})$, and
    \item $C$ and $X$ are unconfounded according to $G(\tilde{M})$,
  \end{enumerate}
a dependence 
  \begin{equation}\label{eq:RCT_CI_statement3}
    X \nIndep_{P_{\tilde{M}}(X,C)} C
  \end{equation}
  implies that $C$ causes $X$ according to $G(\tilde{M})$, and the causal effect of $C$ on $X$ satisfies:
  \begin{equation}\label{eq:RCT_causal_effect3}
    P_{\tilde{M}}\big(X \given \doit(C)\big) = P_{\tilde{M}}(X \given C) \qquad P_{\tilde{M}}(C)\text{-a.s.}.
  \end{equation}
\end{Prp}
\begin{proof}
The first assumption is equivalent to 
$C \hut X \notin G(\tilde{M})$, and the second assumption is equivalent to $C \huh X \notin G(\tilde{M})$.
Hence, out of the eight possible graphs $G(\tilde{M})$, only two satisfy the assumptions:\\
  \begin{center}\scalebox{0.8}{\begin{tikzpicture}
  \begin{scope}
%    \node at (0,-0.8) {$X \Indep C$};
    \node[ndout] (C1) at (-1,0) {$C$};
    \node[ndout] (X1) at (1,0) {$X$};
%    \draw[red, line width=2mm, fill opacity=.30, draw opacity=.30] (-1,-0.5) edge (1,0.5);
%    \draw[red, line width=2mm, fill opacity=.30, draw opacity=.30] (-1,0.5) edge (1,-0.5);
  \end{scope}
  \begin{scope}[xshift=4cm]
%    \node at (0,-0.8) {$C \nIndep X$};
    \node[ndout] (C1) at (-1,0) {$C$};
    \node[ndout] (X1) at (1,0) {$X$};
    \draw[arout] (C1) edge (X1);
  \end{scope}
  \begin{scope}[xshift=8cm]
%    \node at (0,-0.8) {$C \nIndep X$};
    \node[ndout] (C1) at (-1,0) {$C$};
    \node[ndout] (X1) at (1,0) {$X$};
    \draw[arout,bend left] (X1) edge (C1);
    \draw[red, line width=2mm, fill opacity=.30, draw opacity=.30] (-1,-0.5) edge (1,0.5);
    \draw[red, line width=2mm, fill opacity=.30, draw opacity=.30] (-1,0.5) edge (1,-0.5);
  \end{scope}
  \begin{scope}[xshift=12cm]
%    \node at (0,-0.8) {$C \nIndep X$};
    \node[ndout] (C1) at (-1,0) {$C$};
    \node[ndout] (X1) at (1,0) {$X$};
    \draw[arout] (C1) edge (X1);
    \draw[arout,bend left] (X1) edge (C1);
    \draw[red, line width=2mm, fill opacity=.30, draw opacity=.30] (-1,-0.5) edge (1,0.5);
    \draw[red, line width=2mm, fill opacity=.30, draw opacity=.30] (-1,0.5) edge (1,-0.5);
  \end{scope}
  \begin{scope}[yshift=-2cm]
%    \node at (0,-0.8) {$C \nIndep X$};
    \node[ndout] (C1) at (-1,0) {$C$};
    \node[ndout] (X1) at (1,0) {$X$};
    \draw[arlat,bend left] (C1) edge (X1);
    \draw[red, line width=2mm, fill opacity=.30, draw opacity=.30] (-1,-0.5) edge (1,0.5);
    \draw[red, line width=2mm, fill opacity=.30, draw opacity=.30] (-1,0.5) edge (1,-0.5);
  \end{scope}
  \begin{scope}[yshift=-2cm,xshift=4cm]
%    \node at (0,-0.8) {$C \nIndep X$};
    \node[ndout] (C1) at (-1,0) {$C$};
    \node[ndout] (X1) at (1,0) {$X$};
    \draw[arout] (C1) edge (X1);
    \draw[arlat,bend left] (C1) edge (X1);
    \draw[red, line width=2mm, fill opacity=.30, draw opacity=.30] (-1,-0.5) edge (1,0.5);
    \draw[red, line width=2mm, fill opacity=.30, draw opacity=.30] (-1,0.5) edge (1,-0.5);
  \end{scope}
  \begin{scope}[yshift=-2cm,xshift=8cm]
%    \node at (0,-0.8) {$C \nIndep X$};
    \node[ndout] (C1) at (-1,0) {$C$};
    \node[ndout] (X1) at (1,0) {$X$};
    \draw[arout,bend left] (X1) edge (C1);
    \draw[arlat,bend left] (C1) edge (X1);
    \draw[red, line width=2mm, fill opacity=.30, draw opacity=.30] (-1,-0.5) edge (1,0.5);
    \draw[red, line width=2mm, fill opacity=.30, draw opacity=.30] (-1,0.5) edge (1,-0.5);
  \end{scope}
  \begin{scope}[yshift=-2cm,xshift=12cm]
%    \node at (0,-0.8) {$C \nIndep X$};
    \node[ndout] (C1) at (-1,0) {$C$};
    \node[ndout] (X1) at (1,0) {$X$};
    \draw[arout] (C1) edge (X1);
    \draw[arout,bend left] (X1) edge (C1);
    \draw[arlat,bend left] (C1) edge (X1);
    \draw[red, line width=2mm, fill opacity=.30, draw opacity=.30] (-1,-0.5) edge (1,0.5);
    \draw[red, line width=2mm, fill opacity=.30, draw opacity=.30] (-1,0.5) edge (1,-0.5);
  \end{scope}
  \end{tikzpicture}}\end{center}
By the Markov property, if the edge $C \tuh X$ were absent in $G(\tilde{M})$, then $X \Indep_{P_{\tilde{M}}(X,C)} C$. 
In both cases, rule 2 of the causal do-calculus applied to $G(\tilde{M})$ yields the identity \eref{eq:RCT_causal_effect3}. 
\end{proof}

Equation \eref{eq:RCT_CI_statement1} is equivalent to the existence of values $c,c' \in \CX_C$ such that
  $$P_{M}(X \given \doit(C=c)) \ne P_{M}(X \given \doit(C=c')).$$
Equation \eref{eq:RCT_CI_statement3} is equivalent to the existence of values 
  $$P_{\tilde{M}}(X \given C=c) \ne P_{\tilde{M}}(X \given C=c')$$
  for every version of $P_{\tilde{M}}(X \given C)$. 
%(and something similar holds for equation \eref{eq:RCT_CI_statement3}).
These two statements are subtly different. We will see in the next chapter, that as long
as $C$ is discrete, they are actually not that different when testing these statements from 
a finite sample. 

Apart from assuming that there exists a simple iSCM that provides an accurate model,
in both cases, we made the following (implicit or explicit) causal assumptions regarding 
the treatment variable:
\begin{enumerate}
  \item outcome $X$ does not cause treatment $C$ according to $G(M)$;
  \item outcome $X$ and treatment $C$ are unconfounded according to $G(M)$.
\end{enumerate}
As a first remark, we opted here for the strongest version of the notions ``does not cause'' and ``are unconfounded''---the graphical ones.
This might not be \emph{necessary} for deriving the desired conclusions, but it surely is convenient.\footnote{Some researchers prefer potential outcomes because these do not require them to make such graphical assumptions, or even to talk about a graph at all. In our opinion, causal graphs are very convenient to communicate causal modeling assumptions, and we are willing to pay the price for slightly strengthening the untestable assumptions from rung-3 to the graphical rung.}
The first assumption is commonly deemed justified if the outcome is an event that occurs later in time than the treatment event.\footnote{The only exception seems to be certain scenarios involving time travel in science fiction novels/movies.}
The second assumption is usually defended by appealing to randomization.
Indeed, if treatment is decided solely by a coin flip, then it seems hard to imagine some factor that can influence the outcome of the coin flip and also influence the outcome $X$.\footnote{Perhaps an almighty God could be such a factor, though this possibility is often met with skepticism in scientific circles.}
If in addition it is ensured that all subjects are included in the data set (to avoid introducing selection bias), then it is hard to conceive of any remaining source of confounding bias between outcome $X$ and treatment $C$.\footnote{An example in which selection bias would be present is if patients in the placebo group that become too ill drop out of the study (for example, for ethical reasons) and are not included in the data.}

We have shown (in two slightly different ways) that under these assumptions, 
if the distribution of the outcome $X$ differs between the two groups of
patients (``treatment group'' with $C=1$ vs.\ ``control group'' with $C=0$), 
then treatment must be a cause of the outcome, at least in this population of patients.

\begin{comment}
\begin{figure}\centering\begin{tikzpicture}
    \node at (0,2.2) {(a) Two separate data sets:};
    \node[anchor=east] at (-3.5,0) {\scalebox{0.8}{\small\begin{tabular}{c}
      Placebo \\
      ($C=0$): \\
      \hline
      $X$ \\
      \hline
      \color{blue!70}-0.2 \\
      \color{blue!70} 0.6 \\
      \color{blue!70}-1.7 \\
      \color{blue!70}\dots \\
      \hline
    \end{tabular}}};
    \node at (-2,0) {\includegraphics[width=3cm]{RCThistogram0.pdf}};

    \node[anchor=east] at (1.5,0) {\scalebox{0.8}{\small\begin{tabular}{c}
      Drug \\
      ($C=1$): \\
      \hline
      $X$ \\
      \hline
      \color{red!70} -0.3 \\
      \color{red!70} 1.8  \\ 
      \color{red!70} -0.1 \\ 
      \color{red!70} \dots  \\ 
      \hline
    \end{tabular}}};
    \node at (3,0) {\includegraphics[width=3cm]{RCThistogram1.pdf}};

    \begin{scope}[xshift=6.9cm]
    \node at (0,2.2) {(b) Pooled data:};
    \node[anchor=east] at (0,0) {\scalebox{0.8}{\small\begin{tabular}{cc}
      $C$ & $X$ \\
      \hline
      0 & \color{blue!70}-0.2 \\
      0 & \color{blue!70} 0.6 \\
      0 & \color{blue!70}-1.7 \\
      0 & \color{blue!70}\dots \\
      1 & \color{red!70} -0.3 \\
      1 & \color{red!70} 1.8  \\ 
      1 & \color{red!70} -0.1  \\ 
      1 & \color{red!70} \dots \\ 
      \hline
    \end{tabular}}};
%    \node at (0,-5) {\includegraphics[width=3cm]{RCThistogram01.pdf}};
      \node at (1.5,0) {\includegraphics[width=3cm,trim={0cm 3mm 0cm 0cm},clip]{RCThistogram01b.pdf}};
    \end{scope}
  \end{tikzpicture}
  \caption{Illustration of the data from an example randomized controlled trial. When treatment $C$ is
  randomized, the data can either be interpreted
  as (a) two separate data sets, one for the treatment and one for the control group, or (b) 
  as a single data set including a context variable indicating treatment/control. Note that in this
  particular example, $C$ is dependent on $X$ in the pooled data (or equivalently, the distribution of $X$
  differs between contexts $C=0$ and $C=1$), which implies that $C$ is a cause of $X$.\label{fig:RCT_example}}
\end{figure}

Supposing that treatment is completely randomized,
there are two conceptually different ways of testing this in the data, depending on whether
we treat the data as a single pooled data set, or rather as two separate data sets (each one corresponding 
to a particular patient group), see also Figure~\ref{fig:RCT_example}.
To test whether $P(X \given \doit(C=0)) \ne P(X \given \doit(C=1))$,
we can test whether the 
distribution of $X$ is statistically different in the two groups. This can 
be tested with a two-sample test, for example, a $t$-test or a Wilcoxon test. 
The other alternative is to consider the data as a single \emph{pooled} data set. 
The question then becomes whether the
conditional distribution of $X$ given $C=0$ differs from the conditional distribution
of $X$ given $C=1$, i.e., whether $P(X \given C=0) \ne P(X \given C=1)$. 
This can be done with a conditional independence test, for example, Fisher's exact test,
or a partial correlation test.
In the next lecture, we will look in more detail at possible tests.
\end{comment}

If one formally considers treatment as an exogenous input variable, but then also assumes 
that its values are randomly assigned, then the differences between both modeling approaches
are purely cosmetical. However, there is one advantage that the exogenous input approach has over the 
other approach: it does not restrict in any way the possible values of the treatment variable. 
This allows more freedom in the experimental design and sampling scheme design. 
For example, one can decide prior to conducting the RCT that 
the sampling scheme should end up with an equal number of patients in both groups. 
In case treatment is assigned by flipping a coin for each patient, it is rather unlikely
that we will end up with exactly the same number of patients in both groups. 


\subsection{Estimating average treatment effects}

Another way to formalize the randomized controlled setting is by using potential outcomes.
For a binary treatment variable, one introduces two random variables per patient:
$X^{\doit(c_n=1)}_n$ and $X^{\doit(c_n=0)}_n$, corresponding to the potential outcomes for
the $n$'th patient if we treat the patient, or not, respectively. 
Given the actual treatment $C_n$, we then define the actual outcome as 
$X_n := X_n^{\doit(c_n=C_n)}$. In practice, we only observe the actual outcome for
each patient, while the other potential outcome for each patient remains latent. 

An estimand that one is typically interested in and that captures the average strength
of the causal effect of the treatment on the outcome is the so-called
\emph{average treatment effect (ATE)}
$$\tau := \Exp(X_n^{\doit(c_n=1)} - X_n^{\doit(c_n=0)}).$$
To estimate the ATE, one assumes that $C_n$ is randomized. This motivates the assumption
that treatment and outcome are unconfounded:
%\eref{eq:scm_rung3_unconfounded_acyclic}
%$$X_n = X_n^{\doit(C_n)}$$
$$X_n^{\doit(c_n=0)} \Indep C_n, \qquad X_n^{\doit(c_n=1)} \Indep C_n$$
(which follows from Proposition~\ref{prp:scm_rung3_unconfounded_acyclic} if one assumes an underlying iSCM).
One can then show that the difference-in-means estimator 
$$\hat{\tau} := \frac{1}{|\{n : C_n=1\}|} \sum_{\substack{n=1 \\ C_n=1}}^N X_n - \frac{1}{|\{n : C_n=0\}|} \sum_{\substack{n=1 \\ C_n=0}}^N X_n$$
is an unbiased, consistent estimator of the ATE $\tau$.
Curiously enough, while we can speak of the difference
$X_n^{\doit(c_n=1)} - X_n^{\doit(c_n=0)}$ as the individual treatment effect,
this is in many settings an unobservable quantity; however, the \emph{average} treatment effect can 
be estimated from observed data. In the iSCM setting, we can think of the potential
outcomes as counterfactuals in a twin iSCM. However, when assuming an underlying iSCM,
there is no need to go to the counterfactual level, as one can simply define the
ATE as
$$\tau := \Exp_{M}(X \given \doit(C=1)) - \Exp_{M}(X \given \doit(C=0)).$$
In the presence of observed covariates $Z$, one often considers also the \emph{conditional
average treatment effect (CATE)}, which we can define as
$$\Exp_{M}(X \given \doit(C=1),Z) - \Exp_{M}(X \given \doit(C=0),Z).$$
when assuming an underlying iSCM $M$. There is a large body of literature that considers the
question of studying the (asymptotic) efficiency of estimators of the (conditional)
average treatment effect. For a nice account of this surprisingly non-trivial inference problem, 
see e.g.\ \cite{Wager2020}.


\subsection{Local Causal Discovery}
Although the most reliable way to discover causal relations and to estimate their effects
is by means of a randomized controlled trial, it is not always possible or feasible to 
perform such an experiment. One alternative is provided by the Local Causal Discovery
(LCD) algorithm \cite{Cooper1997}.

LCD is a \emph{constraint-based} causal discovery algorithm which means that it discovers
causal relations by combining the results of conditional independence tests on data.
It can be used for the purely observational causal discovery setting where certain background knowledge is available
that is weaker than that for the randomized controlled trial. In particular, no randomization is necessary.

The basic idea behind the LCD algorithm is the following result of \cite{Cooper1997}
(originally formulated for L-CBNs, but easily generalized to simple iSCMs):
\begin{Prp}\label{prp:LCD}
  Let $M$ be a simple SCM with three endogenous variables $V = \{1,2,3\}$ and no exogenous input variables ($J = \emptyset$).
  Suppose that it is $\sigma$-faithful. 
  If $X_2$ is not a cause of $X_1$ according to $G(M)$,
  the following conditional (in)dependencies\footnote{Henceforth, we will no longer always explicitly write
  the Markov kernel as a subscript to the conditional independence symbol if it is clear
  from the context which Markov kernel is meant.} in the observational distribution $P_{M}(X_1,X_2,X_3)$
  $$X_1 \nIndep X_2, \quad X_2 \nIndep X_3, \quad X_1 \Indep X_3 \given X_2$$
  imply that the causal graph $G(M)$ must be one of the three DMGs in Figure~\ref{fig:LCD}. 
  Hence,
  \begin{enumerate}
%  \item $X_2$ is a cause of $X_3$;
  \item $X_3$ is not a cause of $X_2$ according to $G(M)$;
  \item $X_2$ is a direct cause of $X_3$ w.r.t\ $\{1,2,3\}$ according to $G(M)$;
  \item $X_2$ and $X_3$ are unconfounded according to $G(M)$;
  \item the causal effect of $X_2$ on $X_3$ according to $M$ is identified as:
    \begin{equation}\label{eq:LCD}
      P_{M}\big(X_3 \given \doit(X_2)\big) = P_{M}(X_3 \given X_2) \qquad P_M(X_2)\text{-a.s.}.
    \end{equation}
  \end{enumerate}
\end{Prp}
\begin{figure}
  \centering
  \begin{tikzpicture}
    \begin{scope}
      \node[ndout] (X1) at (-3,0) {$X_1$};
      \node[ndout] (X2) at (-1.5,0) {$X_2$};
      \node[ndout] (X3) at (0,0) {$X_3$};
      \draw[arout] (X1) edge (X2);
      \draw[arout] (X2) edge (X3);
    \end{scope}
    \begin{scope}[xshift=5cm]
      \node[ndout] (X1) at (-3,0) {$X_1$};
      \node[ndout] (X2) at (-1.5,0) {$X_2$};
      \node[ndout] (X3) at (0,0) {$X_3$};
      \draw[arlat,bend left] (X1) edge (X2);
      \draw[arout] (X2) edge (X3);
    \end{scope}
    \begin{scope}[xshift=10cm]
      \node[ndout] (X1) at (-3,0) {$X_1$};
      \node[ndout] (X2) at (-1.5,0) {$X_2$};
      \node[ndout] (X3) at (0,0) {$X_3$};
      \draw[arout] (X1) edge (X2);
      \draw[arlat,bend left] (X1) edge (X2);
      \draw[arout] (X2) edge (X3);
    \end{scope}
  \end{tikzpicture}
  \caption{All possible causal graphs detected by LCD.\label{fig:LCD}}
\end{figure}
\begin{proof}
The proof proceeds by enumerating all (possibly cyclic) DMGs on three variables that the causal graph
$G(M)$ could be, and ruling out the ones that do not satisfy the assumptions.
The assumption that $X_2$ is not a cause of $X_1$ according to $G$ implies that there is no directed edge
$X_2 \tuh X_1$ in the graph $G(M)$. If there were an edge between $X_1$ and $X_3$,
$X_1 \Indep X_3 \given X_2$ would not hold ($\sigma$-faithfulness). Also, since $X_1 \nIndep X_2$, $X_1$ and
$X_2$ must be adjacent (Markov property). Similarly, $X_2$ and $X_3$ must be adjacent. $X_2$ cannot be
a collider on any walk between $X_1$ and $X_3$ ($\sigma$-faithfulness). Since the only possible edges between
$X_1$ and $X_2$ are $X_1 \tuh X_2$ and $X_1 \huh X_2$ (both of which are into $X_2$), 
this means that there must be a
directed edge $X_2 \tuh X_3$, but there cannot be a  bidirected edge $X_2 \huh X_3$ or directed edge $X_2 \hut X_3$. 
In other words, the only three possible causal graphs are the ones in Figure~\ref{fig:LCD}.
The causal do-calculus applied to $G(M)$ yields \eref{eq:LCD}.
\end{proof}
Note that one can apply this to the marginalization $M_{\{1,2,3\}}$ in case $V \supseteq \{1,2,3\}$.

In one of the first applications of LCD, it was discovered that nausea causes vomiting 
\cite{SpirtesCooper_AISTATS_99}.
The next example provides another successful application of LCD.
\begin{Eg}\label{eg:lcd_sachs}
  PIP2 and PIP3 are phospholipids that play an important role in human immune system cells.
  Figure~\ref{fig:lcd_sachs} shows a scatter plot of PIP2 and PIP3 expression levels,
  measured in individual human immune system cells, after activation of certain protein signaling cascades in these cells
  \cite{SPP05}.
  The measurements have been performed
  under two different experimental conditions: observational ($C=0$, in blue), and after intervening by administering the chemical
  compound Psitectorigenin to the cells before measuring the PIP2 and PIP3 levels ($C=1$, in red). The experimental protocol
  justifies the assumption that neither PIP2 nor PIP3 expression levels can cause the experimental condition (because
  these expression levels were measured after the experimental condition had been imposed on the cells). 
  Also, assuming the cells to be properly randomized before split into the two groups (corresponding with the
  two experimental conditions) we can conclude that the experimental condition $C$ and PIP2 are unconfounded, 
  and similarly that $C$ and PIP3 are unconfounded. The scatter plot suggests the following conditional independence:
  $$X_{PIP3} \Indep_{P(X_{PIP2},X_{PIP3}\given \doit(X_C))} X_C \given X_{PIP2},$$
  which can also be confirmed with statistical conditional independence tests. Therefore, we have an LCD pattern
  with $X_1=X_C$, $X_2=X_{PIP2}$, $X_3=X_{PIP3}$, which allows us to infer that the PIP2 expression level causes the PIP3 expression level.
  Under the randomization assumption, we can even infer that Psitectorigenin exposure is a cause of PIP2 expression levels.
  This is in line with the Psitectorigenin being known as an inhibitor of PIP2, reducing the quantity of PIP2 in cells after exposure of the cells to this inhibitor.
\begin{figure}\centering
  \includegraphics[width=0.5\textwidth]{LCD_Psitectorigenin_PIP2_PIP3.pdf}
  \caption{``Text-book'' example of an LCD pattern in flow cytometry data of \cite{SPP05}. See Example~\ref{eg:lcd_sachs} for details.\label{fig:lcd_sachs}}
\end{figure}
\end{Eg}
A high-dimensional adaptation has also been shown to be successful in predicting the effects of gene knockout on gene
expression levels from large-scale interventional yeast gene expression data \cite{VersteegMooij_BIBM_19}.

In case more than three variables have been observed, one can run LCD on all triples of variables for which its assumptions apply.
In that case, one should keep in mind that a direct edge in a marginalized graph does not imply the presence of the directed 
edge in the original graph (only the presence of a directed path). In other words, with respect to a larger set of observed variables, 
the causal relations found by LCD are not necessarily direct.

In case of more than three observed variables, one can also replace the single variable $X_2$ in the LCD algorithm by a subset
of variables, a so-called \emph{separating set}. This idea is exploited efficiently in case of many variables in the Invariant Causal Prediction algorithm \cite{ICP2016}.

\begin{comment}
Another possible extension is the following conditional version of LCD. \Joris{TODO}
\begin{Prp}\label{prp:LCD_conditional}
  Let $M$ be a simple iSCM with observed endogenous variables $O = \{1,2,3\} \dot\cup C \subseteq V$ and exogenous input variables $X_J$.
  Suppose that it is $i\sigma$-faithful (w.r.t.\ $O$). 
  If $X_2$ is not a cause of $X_1$ according to $M$,
  the following conditional (in)dependencies in the Markov kernel $P_{M}(X_O \given X_J)$
  $$X_1 \nIndep X_2 \given X_J, X_C \quad \land \quad X_2 \nIndep X_3 \given X_J, X_C \quad \land \quad X_1 \Indep X_3 \given X_2,X_J,X_C$$
  imply that 
  \begin{enumerate}
%  \item $X_2$ is a cause of $X_3$;
  \item $X_3$ is not a cause of $X_2$ according to $M$;
  \item $X_2$ is a direct cause of $X_3$ w.r.t\ $\{1,2,3\}$ according to $M$;
  \item $X_2$ and $X_3$ are unconfounded according to $M$;
  \item the causal effect of $X_2$ on $X_3$ according to $M$ is given by:
    \begin{equation}\label{eq:LCD}
      P_{M}\big(X_3 \given \doit(X_2)\big) = P_{M}(X_3 \given X_2) \qquad P_M(X_2)\text{-a.s.}.
    \end{equation}
  \end{enumerate}
\end{Prp}
\begin{proof}
The proof proceeds by enumerating all (possibly cyclic) DMGs on three variables that the observable graph
$G_{O|J}(M)$ could be, and ruling out the ones that do not satisfy the assumptions.
The assumption that $X_2$ is not a cause of $X_1$ implies that there is no directed edge
$X_2 \tuh X_1$ in the graph $G_{O|J}(M)$. If there were an edge between $X_1$ and $X_3$,
$X_1 \Indep X_3 \given X_2$ would not hold (faithfulness). Also, since $X_1 \nIndep X_2$, $X_1$ and
$X_2$ must be adjacent (Markov property). Similarly, $X_2$ and $X_3$ must be adjacent. $X_2$ cannot be
a collider on any walk between $X_1$ and $X_3$ (faithfulness). Since the only possible edges between
$X_1$ and $X_2$ are $X_1 \tuh X_2$ and $X_1 \huh X_2$ (both of which have an arrowhead at $X_2$), 
this means that there must be a
directed edge $X_2 \tuh X_3$, but there cannot be a  bidirected edge $X_2 \huh X_3$ or directed edge $X_2 \hut X_3$. 
In other words, the only three possible graphs are the ones in Figure~\ref{fig:LCD}.
The causal do-calculus applied to $G_{O|J}(M)$ yields \eref{eq:LCD}.
\end{proof}
\end{comment}

\subsection{Y-structures}

For both the randomized controlled trial and the LCD algorithm, we need prior knowledge: we need to know already that one of the variables is not a cause of another one. 
It turns out that in the absence of any such causal background knowledge, we can sometimes still deduce causal relationships from observed conditional independences.
The simplest such example is given by the ``Y-structure'' pattern \cite{ManiPhD2006}.
We here also give the generalization of the Y-structure pattern to simple SCMs.
\begin{Prp}\label{prp:y-structures}
  Let $M$ be a simple SCM with endogenous variables $V = \{1,2,3,4\}$ and no exogenous input variables ($J = \emptyset$).
  Suppose that it is $\sigma$-faithful. 
  The following conditional (in)dependencies in the observational distribution $P_{M}(X_1,X_2,X_3,X_4)$
  $$\begin{array}{lll}
    X_1 \nIndep X_4,  & X_2 \nIndep X_4, & X_1 \Indep X_2, \\
    X_1 \Indep X_4 \given X_3, & X_2 \Indep X_4 \given X_3, & X_1 \nIndep X_2 \given X_3,
  \end{array}$$
  imply that the causal graph $G(M)$ must be one of the nine DMGs in Figure~\ref{fig:Ystruct}.
  Hence,
  \begin{enumerate}
    \item $X_4$ is not a cause of $X_3$ according to $G(M)$;
    \item $X_3$ is a direct cause of $X_4$ w.r.t.\ $\{1,2,3,4\}$ according to $G(M)$;
    \item $X_3$ and $X_4$ are unconfounded according to $G(M)$;
    \item the causal effect of $X_3$ on $X_4$ according to $M$ satisfies:
  \begin{equation}\label{eq:Y_causal_effect}
    P_{M}\big(X_4 \given \doit(X_3)\big) = P_{M}(X_4 \given X_3) \qquad P_M(X_3)\text{-a.s.}.
  \end{equation}
  \end{enumerate}
\end{Prp}
\begin{proof}
By using the global Markov property and the faithfulness assumption, one
can check that the only (cyclic or acyclic) graphs that are compatible with the observed
conditional independences are the ones in Figure~\ref{fig:Ystruct}.
The statements now follow.
\end{proof}
Note that one can apply this to the marginalization $M_{\{1,2,3,4\}}$ in case $V \supseteq \{1,2,3,4\}$.

This example illustrates how conditional independence patterns in the observational
distribution allow one to infer certain features of the underlying causal model. This
principle is exploited more generally by constraint-based methods, and implicitly,
by score-based methods that optimize a penalized likelihood over (equivalence classes of) causal graphs.
In Chapter~\ref{sec:fci} we will describe in detail one of the most sophisticated constraint-based
causal discovery methods, Fast Causal Inference (FCI).

Typically, the graph cannot be completely identified from purely observational data.
For example, in the Y-structure case, the conditional independences in the observational data
do not allow to conclude whether the dependence between $X_1$ and $X_3$ is explained by
$X_1$ being a cause of $X_3$, or by $X_1$ and $X_3$ having a latent common cause,
or because of some latent selection mechanism, or a combination of those.
However, under an appropriate faithfulness assumption, one can deduce the Markov
equivalence class of the graph from the conditional independences in the observational data, 
i.e., the class of all DMGs that induce the same separations.

A significant practical disadvantage of causal discovery methods from purely observational data is that
they typically need \emph{very large} sample sizes and \emph{strong assumptions} in order to work reliably (even on simulations).

\begin{figure}\centering
  \begin{tikzpicture}
    \begin{scope}
      \node[ndout] (X1) at (-1,0) {$X_1$};
      \node[ndout] (X2) at (1,0) {$X_2$};
      \node[ndout] (X3) at (0,-1) {$X_3$};
      \node[ndout] (X4) at (0,-2.5) {$X_4$};
      \draw[arout] (X1) edge (X3);
  %    \draw[arlat,bend right] (X1) edge (X3);
      \draw[arout] (X2) edge (X3);
  %    \draw[arlat,bend left] (X2) edge (X3);
      \draw[arout] (X3) edge (X4);
    \end{scope}
    \begin{scope}[xshift=5cm]
      \node[ndout] (X1) at (-1,0) {$X_1$};
      \node[ndout] (X2) at (1,0) {$X_2$};
      \node[ndout] (X3) at (0,-1) {$X_3$};
      \node[ndout] (X4) at (0,-2.5) {$X_4$};
  %    \draw[arout] (X1) edge (X3);
      \draw[arlat,bend right] (X1) edge (X3);
      \draw[arout] (X2) edge (X3);
  %    \draw[arlat,bend left] (X2) edge (X3);
      \draw[arout] (X3) edge (X4);
    \end{scope}
    \begin{scope}[xshift=10cm]
      \node[ndout] (X1) at (-1,0) {$X_1$};
      \node[ndout] (X2) at (1,0) {$X_2$};
      \node[ndout] (X3) at (0,-1) {$X_3$};
      \node[ndout] (X4) at (0,-2.5) {$X_4$};
      \draw[arout] (X1) edge (X3);
      \draw[arlat,bend right] (X1) edge (X3);
      \draw[arout] (X2) edge (X3);
  %    \draw[arlat,bend left] (X2) edge (X3);
      \draw[arout] (X3) edge (X4);
    \end{scope}
    \begin{scope}[yshift=-4cm]
      \node[ndout] (X1) at (-1,0) {$X_1$};
      \node[ndout] (X2) at (1,0) {$X_2$};
      \node[ndout] (X3) at (0,-1) {$X_3$};
      \node[ndout] (X4) at (0,-2.5) {$X_4$};
      \draw[arout] (X1) edge (X3);
  %    \draw[arlat,bend right] (X1) edge (X3);
  %    \draw[arout] (X2) edge (X3);
      \draw[arlat,bend left] (X2) edge (X3);
      \draw[arout] (X3) edge (X4);
    \end{scope}
    \begin{scope}[xshift=5cm,yshift=-4cm]
      \node[ndout] (X1) at (-1,0) {$X_1$};
      \node[ndout] (X2) at (1,0) {$X_2$};
      \node[ndout] (X3) at (0,-1) {$X_3$};
      \node[ndout] (X4) at (0,-2.5) {$X_4$};
  %    \draw[arout] (X1) edge (X3);
      \draw[arlat,bend right] (X1) edge (X3);
  %    \draw[arout] (X2) edge (X3);
      \draw[arlat,bend left] (X2) edge (X3);
      \draw[arout] (X3) edge (X4);
    \end{scope}
    \begin{scope}[xshift=10cm,yshift=-4cm]
      \node[ndout] (X1) at (-1,0) {$X_1$};
      \node[ndout] (X2) at (1,0) {$X_2$};
      \node[ndout] (X3) at (0,-1) {$X_3$};
      \node[ndout] (X4) at (0,-2.5) {$X_4$};
      \draw[arout] (X1) edge (X3);
      \draw[arlat,bend right] (X1) edge (X3);
  %    \draw[arout] (X2) edge (X3);
      \draw[arlat,bend left] (X2) edge (X3);
      \draw[arout] (X3) edge (X4);
    \end{scope}
    \begin{scope}[yshift=-8cm]
      \node[ndout] (X1) at (-1,0) {$X_1$};
      \node[ndout] (X2) at (1,0) {$X_2$};
      \node[ndout] (X3) at (0,-1) {$X_3$};
      \node[ndout] (X4) at (0,-2.5) {$X_4$};
      \draw[arout] (X1) edge (X3);
  %    \draw[arlat,bend right] (X1) edge (X3);
      \draw[arout] (X2) edge (X3);
      \draw[arlat,bend left] (X2) edge (X3);
      \draw[arout] (X3) edge (X4);
    \end{scope}
    \begin{scope}[xshift=5cm,yshift=-8cm]
      \node[ndout] (X1) at (-1,0) {$X_1$};
      \node[ndout] (X2) at (1,0) {$X_2$};
      \node[ndout] (X3) at (0,-1) {$X_3$};
      \node[ndout] (X4) at (0,-2.5) {$X_4$};
  %    \draw[arout] (X1) edge (X3);
      \draw[arlat,bend right] (X1) edge (X3);
      \draw[arout] (X2) edge (X3);
      \draw[arlat,bend left] (X2) edge (X3);
      \draw[arout] (X3) edge (X4);
    \end{scope}
    \begin{scope}[xshift=10cm,yshift=-8cm]
      \node[ndout] (X1) at (-1,0) {$X_1$};
      \node[ndout] (X2) at (1,0) {$X_2$};
      \node[ndout] (X3) at (0,-1) {$X_3$};
      \node[ndout] (X4) at (0,-2.5) {$X_4$};
      \draw[arout] (X1) edge (X3);
      \draw[arlat,bend right] (X1) edge (X3);
      \draw[arout] (X2) edge (X3);
      \draw[arlat,bend left] (X2) edge (X3);
      \draw[arout] (X3) edge (X4);
    \end{scope}
  \end{tikzpicture}\hfill
%  \includegraphics[width=0.26\textwidth]{Ystruct_12.pdf}
%  \includegraphics[width=0.26\textwidth]{Ystruct_14.pdf}
%  \includegraphics[width=0.26\textwidth]{Ystruct_24.pdf}
  \caption{Causal graphs satisfying the ``Y-structure'' pattern on four variables.
%Right: Scatter plots illustrating the Y-structure pattern in purely observational data, where $X_3$ is discrete-valued and its value is indicated by color (red/blue).
  \label{fig:Ystruct}}
\end{figure}

